\ansclass{2}{Estimation \& Performance Math}
\ans{1} (a) Divide by $10^5$: \final{$\approx$ 10/s, 1\,000/s and 10\,000/s} (exactly 11.6, 1\,157 and 11\,574: the rule of thumb is 15\% low, fine for a decision). (b) $2 \times 10^6 \times 20 = 4 \times 10^7$ reads/day, $/10^5 =$ \final{400 reads/s on average, $\approx$ 1\,200/s at the peak}. One well-tuned app server; a cache makes it trivial. (c) \final{1.25 GB/s} ($10/8$). (d) $1\,200 \times 1$ MB $= 1.2$ GB/s $=$ \final{9.6 Gbit/s $\approx$ 10 Gbit/s}, a whole network card; with 95\% served by the CDN the origin sends \final{$\approx$ 0.5 Gbit/s}.
\careful{mixing bits and bytes: network links are sold in bits per second, files are measured in bytes. A factor of 8 is the difference between ``fits'' and ``does not fit''.}
\ans{2} (a) $60\,000 \times 0.5$ ms $=$ \final{30 s}, only in round trips. (b) $60\,000 \times 50$ B $= 3$ MB; on 10 Gbit/s (1.25 GB/s) that is 2.4 ms. Total \final{$\approx$ 20--40 ms}, dominated by the database work. (c) 300 KB take 0.24 ms on the wire: \final{$\approx$ 1--2 ms}. The three versions are four orders of magnitude apart. (d) $50 \times 0.2$--$0.5$ ms $=$ \final{10--25 ms} of pure waiting; \final{batch with \texttt{MGET} or pipelining} (one round trip).
\careful{counting only the work of each call and forgetting that sequential calls pay one round trip each.}
\ans{3} (a) $L = \lambda W = 1\,000 \times 0.040 =$ \final{40 holds in flight}, so at least \final{10 threads per replica} (plus headroom). (b) $1\,000 \times 0.010 =$ \final{10 connections busy}; a pool of $\approx$ 15 across all replicas is enough. \final{200 connections add no capacity}: the database still has 8 slots, and the extra connections only queue inside it (and cost memory). (c) $\rho = \lambda S / c = 1\,000 \times 0.005 / 8 =$ \final{62.5\%}; $X_{\max} = c/S = 8/0.005 =$ \final{1\,600 holds/s}. (d) $100 \times 0.3 = 30$ in flight before; $100 \times 3 =$ \final{300 after, more than the 200 threads}: \texttt{order} runs out of threads and starts failing even requests that do not need the provider. Fix: a timeout and a bulkhead (class 10).
\careful{reading Little's Law as a statement about CPU. A slow dependency fills the pool with \emph{waiting} threads while the CPU is idle.}
\ans{4} (a) $R$: \final{10, 16.7, 50 and 100 ms}; p99: \final{46, 77, 230 and 461 ms}. From 50\% to 90\% (less than twice the load) the mean gets 5$\times$ worse. (b) $1/(1-\rho)$ goes from $1/0.3 = 3.3$ to $1/0.08 = 12.5$: \final{3.75$\times$}. The model with $S = 5$ ms gives a p99 of $4.6 \times 5/0.08 \approx$ \final{290 ms}, what was observed (and 3.75 $\times$ the 77 ms at 70\%). With a 300 ms target, holds start to time out, clients retry, the retries add load, $\rho$ goes to 1 and the error rate to 100\%. (c) $4 \times 60\% / 3 =$ \final{80\%}: survivable. From 80\%: $4 \times 80\%/3 =$ \final{107\%}, overload as soon as a replica dies.
\careful{thinking that latency grows in proportion to load. It grows like $1/(1-\rho)$: slowly, then all at once.}
\ans{5} (a) $1 - 0.99^N$: \final{4.9\%, 18\% and 63\%}. With 100 calls, most pages are slow although each backend is ``fast at p99''. (b) $1 - 0.99^{20} \approx$ \final{18\%} of searches wait at least 120 ms (the user-facing p82 is the shard p99); with 10 shards, \final{$\approx$ 9.6\%}. (c) \final{Percentiles cannot be averaged.} The fleet p99 depends on how many requests each pod served; merge the histograms (sum the buckets) and compute the percentile once. (d) \final{Hedging}: after the p95, send a second copy to another replica and keep the first answer (about 5\% more requests). Never for a non-idempotent request such as \texttt{POST /orders}. Also: partial results after a deadline.
\careful{reporting the average latency of each backend and concluding that the page is fast.}
\ans{6} (a) $0.9999 \times 0.9995^4 \times 0.999 \times 0.995 =$ \final{99.19\%}, $\approx$ \final{5.8 h down per 30 days}. (b) Without the e-mail: \final{99.69\%} ($\approx$ 2.2 h). (c) $1 - 0.001^2$ for the payment link: \final{99.79\%} ($\approx$ 1.5 h). (d) The four internal components at 99.95\% still cost 0.2\%. Raise each of them to 99.99\%: $0.9999^5 \times (1 - 0.001^2) =$ \final{99.95\%} ($\approx$ 22 min). The rule of the extra nine: \final{in a chain of about five links, each link needs one nine more than the target.} (e) $720/722 = 99.72\%$; halving the failures: $1\,440/1\,442 = 99.86\%$; replacing in 2 min: $720/720.033 =$ \final{99.995\%}. Faster recovery is the bigger and usually cheaper win.
\careful{adding redundancy to the strong links and forgetting the weak, synchronous one (the e-mail).}
\ans{7} (a) $2 \times 10^6 / 60 \approx$ \final{33\,000 hold attempts/s} (more at the edge, with refreshes). (b) $8/0.005 = 1\,600$/s at 100\%, \final{$\approx$ 1\,100/s at 70\%}: demand is \final{20 to 30 times} the capacity. No ``faster Postgres'' closes that gap. (c) 60\,000 seats need 30\,000 fans: \final{$\approx$ 30 s} at 1\,000/s. The last of 2 M fans would wait $2 \times 10^6 / 1\,000 = 2\,000$ s $\approx$ \final{33 min}, but admission stops once the seats are gone (the others see ``sold out''). (d) Orders: $10^5 \times 1$ KB $= 100$ MB/day $\approx$ 37 GB/year; 5 years $\times$ 3 replicas $\approx$ \final{550 GB: no sharding}. PDFs: 20 GB/day $\approx$ 7.3 TB/year, \final{$\approx$ 36 TB in 5 years}: object store, or generate them on demand. (e) ``One PostgreSQL primary handles about 1\,100 holds/s at a safe 70\%, while the on-sale asks for 33\,000 per second. We keep one database and put a waiting room in front of it that admits 1\,000 fans per second; the 60\,000 seats go in about 30 seconds.''
\careful{answering the question ``is Postgres fast?'' instead of comparing demand with capacity at a safe utilization.}
